\section{Source audit and proposed integration}
\label{audit:section}

\subsection{What the audit established}

No new false theorem was found in the specific source arguments used
for the present results. In particular, the incoming cubic bridge,
the centered parity criterion, and the previously proved
trigamma-square moments survive the checks relevant to their use
here. Their stated limits are research questions, not flaws in
their proofs. This is a targeted mathematical audit of the relevant
material, not a certification of every page in the repository.

The appropriate source changes are therefore progress updates and
additional normalization examples. Table~\ref{audit:status-table}
gives the precise mapping. The supplied integration notes identify
the archive-internal source paths and theorem labels; no source
repository was modified in preparing this package.

\begin{table}[htbp]
\centering\small
\begin{tabular}{@{}>{\raggedright\arraybackslash}p{.28\linewidth}>{\raggedright\arraybackslash}p{.36\linewidth}>{\raggedright\arraybackslash}p{.28\linewidth}@{}}
\toprule
Source question & Result established here & Remaining scope\\
\midrule
Cubic harmonic report, Q6 & Quartic power-to-harmonic reduction,
Theorem~\ref{q4:main} & Powers $k\ge5$ and arithmetic reduction of
the retained depth-three coefficient\\[5pt]
Harmonic-parity report, R1 & All pure two-even-tail moments and their
generator, Theorems~\ref{tail:all-moments}, \ref{tail:generator}
& Additional harmonic factors and products of four or more tails\\[5pt]
Harmonic-parity report, R4 & Complex powers, joint resonant completion,
all logarithmic moments, Theorem~\ref{power:completion}
& Products of distinct resolvents or nonlinear continuous-order factors\\[5pt]
Higher-reflection report, centered products & First mixed pair and all
two-tail orders & Higher-depth products and spectral derivatives\\[5pt]
Canonical Gaussian $S_6$, revised $S_8$ & No proof or refutation obtained
& Exact functional bridge to the frozen candidates\\
\bottomrule
\end{tabular}
\caption{Scoped updates relative to the pinned source revision.}
\label{audit:status-table}
\end{table}

\subsection{A concrete correction to a tempting extension}

There is a mathematically incorrect extension that the new results
rule out: replacing a specialized coordinate finite part by the
meromorphic transform at every point where the specialized integrand
happens to be integrable. The counterexample at the end of
Section~\ref{power:sec} is exact. With $p=2$, $u=1$, and
$\lambda=1$, the ordinary integral of
$\partial_x^2\zeta(0,x)R_z(x)$ equals zero. The fixed-$\lambda$
meromorphic continuation of the spectral formula equals $-1$.
The local monomial has a coefficient that vanishes at the same
point as its integration denominator, leaving the nonzero limit
$u-2$.

The correction is to retain the entire moving amplitude before
coefficient extraction, as in \eqref{power:H}. Removing just the
negative Laurent coefficients of a one-variable expansion is
insufficient. At negative integer powers the two smooth endpoint
residues cancel, yet their quotient has the nonzero limit
\eqref{power:smoothcontact}. A cancellation of residues therefore
does not imply that there is no finite endpoint correction.

This is a warning and an exact regression example for future
extensions. The inspected incoming resolvent theorem states its
own convergence and completion conditions; this article does not
attribute the incorrect extension to that theorem. The package
includes a deliberately corrupted coefficient check so that
omitting the endpoint phase is visibly rejected by exact algebra.

\subsection{Four conventions that should accompany integration}

\begin{enumerate}
\item \textbf{Keep strict harmonic coordinates named and normalized.}
The quartic coefficient $\theta$ is the coefficient of $u$ in
$Z_3(1+u,1,1)$, not the constant coefficient and not an order
derivative of a regular function before its poles have been
specified. The entire principal part and constant coefficient
are given in \eqref{q4:theta}. The retained coefficient is
genuine remaining arithmetic data in the present reduction;
its integral representation is a normalization, not an ordinary-
constant evaluation.

\item \textbf{Retain subtractions inside convergent summands.}
The sums in \eqref{tail:ordinary-formula} and
\eqref{tail:rational-polygamma} are ordinary convergent sums of
the displayed brackets. Their polynomial pieces must not be
separated and summed as ordinary series. The equality with the
continued Dirichlet series is proved using the half-shift and
Hurwitz trivial zeros.

\item \textbf{Use the regularized Gauss function at parameter resonances.}
The definition \eqref{moment:eq:regularized} remains valid when
its third parameter is a nonpositive integer. A quotient written
informally as ${}_2F_1/\Gamma(c)$ must be read through that
definition. Likewise the logarithmic Gamma-quotient series has
the smaller disk stated in \eqref{moment:eq:hgenerator}; zeros
of the quotient obstruct its logarithm before they obstruct
the quotient itself.

\item \textbf{Distinguish analytic proofs, finite algebra, and numerical evidence.}
The infinite identities follow from the proofs. Symbolic checks
verify selected finite polynomial consequences. The numerical
files use ordinary multiprecision arithmetic with specified
analytic tail estimates. They are not interval-enclosed evaluations
and do not certify arithmetic independence of any constants.
\end{enumerate}

\subsection{The Gaussian conjectures and the exact obstruction}

The frozen $S_6$ target is in
\texttt{chapters/04-cyclotomic-quotients.tex}; the revised $S_8$
candidate is in \texttt{chapters/04-S8-candidate.tex} of the
manuscript at the pinned revision~\cite{ProveIt}. Both remain
conjectures. Their existing short candidate formulas have not
been replaced by a numerical fit in this work.

The exact-resonant report already proves a specific obstruction:
the $S_6$ target is separated from the entire convergent Cayley
shuffle ideal together with a specified family of $5{,}131$
additional relations~\cite{IncomingExact}. This is stronger than
an unsuccessful small linear solve. It is also a statement about
that named formal relation system, not a proof that the numerical
period conjecture is false or that no other functional relation
can prove it.

The larger search in the independent-orders report generated
$51{,}244$ double-shuffle rows in $24{,}514$ imaginary word
coordinates. Its modular elimination stopped after $8{,}500$
pivots, before exhausting the system~\cite{IncomingIndependent}.
The incomplete search cannot be cited as an exhausted obstruction,
and it supplies no rational identity certificate. No new claim
about its terminal rank is made here.

The new complex-power transform provides analytic identities that
can be tested against a Gaussian target, but this article has not
proved a specialization producing either frozen residual. The
unresolved step is an explicit functional bridge, followed by
exact coefficient extraction. Giving another representation of
an unrelated zeta constant does not accomplish that step.

